% problem_id: S001-lemma-algebraic-center
% source_id: S001 (Stacks Project)
% source_file: groupoids.tex
% statement_locator: groupoids.tex lines 1482--1486
% proof_locator: groupoids.tex lines 1488--1579
% env: lemma
% pairing: tight (verified)
% status: complete (statement + author proof verbatim + reason + locators)
%
\begin{lemma}
\label{lemma-algebraic-center}
Let $k$ be a field. Let $G$ be a locally algebraic group scheme over $k$.
Then the center of $G$ is a closed subgroup scheme of $G$.
\end{lemma}

\begin{proof}
Let $\text{Aut}(G)$ denote the contravariant functor on the category of
schemes over $k$ which associates to $S/k$ the set of automorphisms
of the base change $G_S$ as a group scheme over $S$. There is a natural
transformation
$$
G \longrightarrow \text{Aut}(G),\quad
g \longmapsto \text{inn}_g
$$
sending an $S$-valued point $g$ of $G$ to the inner automorphism of $G$
determined by $g$. The center $C$ of $G$ is by definition the kernel of
this transformation, i.e., the functor which to $S$ associates those
$g \in G(S)$ whose associated inner automorphism is trivial. The statement
of the lemma is that this functor is representable by a closed subgroup
scheme of $G$.

\medskip\noindent
Choose an integer $n \geq 1$. Let $G_n \subset G$ be the $n$th infinitesimal
neighbourhood of the identity element $e$ of $G$. For every scheme $S/k$
the base change $G_{n, S}$ is the $n$th infinitesimal neighbourhood of
$e_S : S \to G_S$. Thus we see that there is a natural transformation
$\text{Aut}(G) \to \text{Aut}(G_n)$ where the right hand side is the
functor of automorphisms of $G_n$ as a scheme ($G_n$ isn't in general
a group scheme). Observe that $G_n$ is the spectrum of an artinian
local ring $A_n$ with residue field $k$ which has finite dimension
as a $k$-vector space
(Varieties, Lemma \ref{varieties-lemma-algebraic-scheme-dim-0}).
Since every automorphism of $G_n$ induces in particular an invertible
linear map $A_n \to A_n$, we obtain transformations of functors
$$
G \to \text{Aut}(G) \to \text{Aut}(G_n) \to \text{GL}(A_n)
$$
The final group valued functor is representable, see
Example \ref{example-general-linear-group}, and the
last arrow is visibly injective.
Thus for every $n$ we obtain a closed subgroup scheme
$$
H_n = \Ker(G \to \text{Aut}(G_n)) = \Ker(G \to \text{GL}(A_n)).
$$
As a first approximation we set $H = \bigcap_{n \geq 1} H_n$
(scheme theoretic intersection). This is a closed subgroup scheme
which contains the center $C$.

\medskip\noindent
Let $h$ be an $S$-valued point of $H$ with $S$ locally Noetherian.
Then the automorphism $\text{inn}_h$ induces the identity on all
the closed subschemes $G_{n, S}$. Consider the kernel
$K = \Ker(\text{inn}_h : G_S \to G_S)$.
This is a closed subgroup scheme of $G_S$ over $S$ containing the
closed subschemes $G_{n, S}$ for $n \geq 1$.
This implies that $K$ contains an open neighbourhood of
$e(S) \subset G_S$, see
Algebra, Remark \ref{algebra-remark-intersection-powers-ideal}.
Let $G^0 \subset G$ be as in Proposition \ref{proposition-connected-component}.
Since $G^0$ is geometrically irreducible, we conclude that $K$ contains
$G^0_S$ (for any nonempty open $U \subset G^0_{k'}$ and any field extension
$k'/k$ we have $U \cdot U^{-1} = G^0_{k'}$, see proof of
Lemma \ref{lemma-irreducible-group-scheme-over-field-qc}).
Applying this with $S = H$ we find that $G^0$ and $H$
are subgroup schemes of $G$ whose points commute: for any scheme $S$
and any $S$-valued points $g \in G^0(S)$, $h \in H(S)$ we have
$gh = hg$ in $G(S)$.

\medskip\noindent
Assume that $k$ is algebraically closed. Then we can pick a $k$-valued
point $g_i$ in each irreducible component $G_i$ of $G$. Observe that in
this case the connected components of $G$ are the irreducible components
of $G$ are the translates of $G^0$ by our $g_i$. We claim that
$$
C = H \cap \bigcap\nolimits_i \Ker(\text{inn}_{g_i} : G \to G)
\quad (\text{scheme theoretic intersection})
$$
Namely, $C$ is contained in the right hand side. On the other hand, every
$S$-valued point $h$ of the right hand side commutes with $G^0$
and with $g_i$ hence with everything in $G = \bigcup G^0g_i$.

\medskip\noindent
The case of a general base field $k$ follows from the result for the
algebraic closure $\overline{k}$ by descent. Namely, let
$A \subset G_{\overline{k}}$ the closed subgroup scheme representing
the center of $G_{\overline{k}}$. Then we have
$$
A \times_{\Spec(k)} \Spec(\overline{k}) =
\Spec(\overline{k}) \times_{\Spec(k)} A
$$
as closed subschemes of $G_{\overline{k} \otimes_k \overline{k}}$ by
the functorial nature of the center. Hence we see that $A$ descends
to a closed subgroup scheme $Z \subset G$ by
Descent, Lemma \ref{descent-lemma-closed-immersion}
(and Descent, Lemma \ref{descent-lemma-descending-property-closed-immersion}).
Then $Z$ represents $C$ (small argument omitted) and the proof is complete.
\end{proof}
